\section{Constante gravitatoria intr\'inseca y representaci\'on tensorial}
\label{rm:ch:gravity}

\subsection{Cadena constitutiva del sector gravitatorio: escala cronológica \texorpdfstring{\(108\)}{108}, reducción tensorial, holonomía y evaluación metrológica de \texorpdfstring{\(G\)}{G}}

La gravedad se organiza aqu\'i en cuatro estratos que no deben fundirse:

\begin{equation}
\label{rm:eq:gravity-causal-chain}
\begin{aligned}
\text{estructura cronol\'ogica CT108 y rectas de acci\'on}
&\longrightarrow \text{selector de periodicidad }\theta_{108}^{\rm circ}
\longrightarrow G_{108},\\
\text{incidencia icosa\'edrica}
&\longrightarrow V_{12}\xrightarrow{P_5}V_5
\xrightarrow{\Gamma_5}\operatorname{STF}_3
\xrightarrow{\Lambda(n)}\mathcal H_{\rm TT}(n),\\
\text{ruta holon\'omica enriquecida}
&\dashrightarrow G_{\rm hol},\\
\text{carta metrol\'ogica}
&\longrightarrow 6.67430\times10^{-11}\,
\mathrm{m^3\,kg^{-1}\,s^{-2}}.
\end{aligned}
\end{equation}

La primera l\'inea expresa una condici\'on intr\'inseca y exacta de
periodicidad. La segunda construye
el portador y su reducci\'on transversal sin traza. La tercera conserva una
frontera de construcci\'on expl\'icita. La cuarta es una evaluaci\'on en una
carta decimal declarada: identifica la magnitud, pero no selecciona el car\'acter
gravitatorio.

\subsection{Recta de magnitud gravitatoria y escala cronol\'ogica \texorpdfstring{\(108\)}{108}}

Sean \(c_{\rm int}>0\), \(t_0>0\),
\(\ell_0=c_{\rm int}t_0\) y una secci\'on positiva
\(\hbar_a\) de la recta de acci\'on, con
\(a\in\{\mathrm{pre},\mathrm{ret}\}\). El transductor dimensional es

\begin{equation}
\label{rm:eq:gravity-transducer}
G_a(\theta)
=\theta\frac{c_{\rm int}^{5}t_0^2}{\hbar_a}
=\theta\frac{c_{\rm int}^{3}\ell_0^2}{\hbar_a},
\qquad \theta>0.
\end{equation}

La igualdad de las dos expresiones usa s\'olo
\(\ell_0=c_{\rm int}t_0\), y
\([G]=L^3M^{-1}T^{-2}\). La dimensi\'on queda fijada por la recta; el
problema matem\'atico es seleccionar el car\'acter adimensional \(\theta\).

El calendario CT108 define

\begin{equation}
\label{rm:eq:gravity-clock108}
T_{108}=108t_0,
\qquad
\omega_{108}=\frac{2\pi}{108t_0},
\qquad
R_{108}=\frac{c_{\rm int}}{\omega_{108}}
=\frac{54}{\pi}\ell_0.
\end{equation}

Para cada hoja de acci\'on, sean

\begin{equation}
\label{rm:eq:gravity-clock-energy-mass}
E_{108,a}=\hbar_a\omega_{108},
\qquad
m_{108,a}=\frac{E_{108,a}}{c_{\rm int}^2}.
\end{equation}

Entonces, la longitud de Compton reducida coincide con el radio del ciclo:

\begin{equation}
\label{rm:eq:gravity-compton-clock}
\bar\lambda_{C,a}
=\frac{\hbar_a}{m_{108,a}c_{\rm int}}
=R_{108}.
\end{equation}

\begin{theorem}[Unicidad del selector bajo la condici\'on de periodicidad HMT]
\label{rm:thm:gravity-circular-selector}
Con la convenci\'on reducida
\(r_g=Gm/c_{\rm int}^2\) y la premisa de periodicidad HMT que identifica el radio
gravitatorio del cuanto CT108 con su radio de fase, la condici\'on de cierre
\begin{equation}
\label{rm:eq:gravity-closing-condition}
r_{g,a}(\theta)=R_{108}
\end{equation}
selecciona una \`unica soluci\'on positiva,
\begin{equation}
\label{rm:eq:gravity-theta108}
\boxed{
\theta_{108}^{\rm circ}=\left(\frac{54}{\pi}\right)^2
=295.4525715010569415303525\ldots .}
\end{equation}
En esa realizaci\'on,
\begin{equation}
\label{rm:eq:gravity-four-lengths}
R_{108}=\bar\lambda_{C,a}=r_{g,a}=\ell_P^{\rm circ}.
\end{equation}
\end{theorem}

\begin{proof}
Por \cref{rm:eq:gravity-transducer,rm:eq:gravity-clock-energy-mass},
\[
r_{g,a}(\theta)
=\frac{G_a(\theta)m_{108,a}}{c_{\rm int}^2}
=\theta c_{\rm int}t_0^2\omega_{108}
=\theta\frac{\pi}{54}\ell_0.
\]
Al compararlo con
\(R_{108}=(54/\pi)\ell_0\) se obtiene necesariamente
\(\theta=(54/\pi)^2\). La ecuaci\'on es lineal en \(\theta\) y todos los
factores son positivos, luego la soluci\'on es \`unica.
\end{proof}

La igualdad \(r_g=R_{108}\) no se deduce de la dimensionalidad de
\cref{rm:eq:gravity-transducer}. Es la condición geométrica explícita de esta
realizaci\'on HMT: acci\'on, Compton y retorno peri\'odico deben determinar una misma
longitud. El teorema demuestra que, aceptada esa condición estructural, el
carácter \(\theta\) queda fijado sin usar el valor SI de \(G\).

\begin{remark}[Convenci\'on de radio]
El teorema usa el radio gravitatorio reducido \(Gm/c^2\). Si se sustituye
por el radio de Schwarzschild \(2Gm/c^2\), el selector cambia por un factor
dos. Las dos expresiones corresponden a condiciones geom\'etricas distintas,
por lo que la convenci\'on debe declararse antes del c\'alculo.
\end{remark}

\subsection{Sistema de Planck y covariancia bidireccional}

\begin{corollary}[Familia de Planck asociada a la estructura cronol\'ogica CT108]
\label{rm:cor:gravity-planck-family}
Para \(G_a=G_a(\theta_{108}^{\rm circ})\),
\begin{align}
\ell_P&=\frac{54}{\pi}\ell_0,
&t_P&=\frac{54}{\pi}t_0,
\label{rm:eq:gravity-planck-length-time}\\
E_{P,a}&=\frac{\hbar_a}{t_P}
=\frac{\pi\hbar_a}{54t_0}
=\frac{h_a}{108t_0},
&\ell_P^2&=\frac{G_a\hbar_a}{c_{\rm int}^3}
=\theta_{108}^{\rm circ}\ell_0^2.
\label{rm:eq:gravity-planck-energy-area}
\end{align}
Adem\'as,
\begin{equation}
\label{rm:eq:gravity-planck-force-power-density}
F_{P,a}=\frac{c_{\rm int}^4}{G_a},
\qquad
P_{P,a}=\frac{c_{\rm int}^5}{G_a},
\qquad
\rho_{P,a}=\frac{\hbar_a}
{(\theta_{108}^{\rm circ})^2c_{\rm int}^5t_0^4}.
\end{equation}
\end{corollary}

\begin{proof}
Las tres primeras identidades proceden de
\(\ell_P=R_{108}\), \(t_P=\ell_P/c_{\rm int}\) y
\(h_a=2\pi\hbar_a\). Para el \`area,
\[
\frac{G_a\hbar_a}{c_{\rm int}^3}
=\theta_{108}^{\rm circ}c_{\rm int}^2t_0^2
=\theta_{108}^{\rm circ}\ell_0^2.
\]
Las expresiones restantes son las composiciones dimensionales de fuerza,
potencia y masa por volumen con la misma secci\'on.
\end{proof}

Sea
\begin{equation}
\label{rm:eq:gravity-Ract}
R_{\rm act}:=\frac{\hbar_{\rm pre}}{\hbar_{\rm ret}}.
\end{equation}
Como la acci\'on aparece en el numerador de la masa cronol\'ogica y en el
denominador de \(G\),
\begin{equation}
\label{rm:eq:gravity-twoway-covariance}
\frac{m_{108,\rm pre}}{m_{108,\rm ret}}=R_{\rm act},
\qquad
\frac{G_{\rm pre}}{G_{\rm ret}}=R_{\rm act}^{-1},
\qquad
G_a\hbar_a
=\theta_{108}^{\rm circ}c_{\rm int}^5t_0^2.
\end{equation}
Por tanto, \(\ell_P^2=G_a\hbar_a/c^3\) no cambia de hoja. La gravedad y
la acci\'on se orientan rec\'iprocamente, mientras el \`area conserva la
geometr\'ia com\'un. Esta cancelaci\'on es la interfaz de confluencia hacia el operador de
\`area del \cref{rm:ch:modular}.

\subsection{Componente icosaédrica de dimensión \(5\)}

El antecedente no es \(\R\) ni una elecci\'on posterior de cinco
polarizaciones. El objeto encargado de transportar la incidencia se define
expl\'icitamente antes de introducir la realizaci\'on excepcional.

Sea \(X_{\rm TPK}^{\rm exc}\) el subespacio del estado \(\TPK\) enriquecido
que conserva, adem\'as de la salida aritm\'etica, la orientaci\'on, la ruta,
el acarreo y la incidencia excepcional. Para explicitar el espacio geom\'etrico de llegada,
sea \(\varphi_g=(1+\sqrt5)/2\), \(\nu=\sqrt{1+\varphi_g^2}\), y definamos
\begin{equation}
\label{rm:eq:gravity-icosahedron-vertices}
\Omega_{12}^{\rm ico}
=\frac1\nu\bigl(
\{0\}\times\{\pm1\}\times\{\pm\varphi_g\}
\;\cup\;
\{\pm1\}\times\{\pm\varphi_g\}\times\{0\}
\;\cup\;
\{\pm\varphi_g\}\times\{0\}\times\{\pm1\}
\bigr)\subset S^2.
\end{equation}
Sus doce elementos son las clases orientadas de v\'ertice de un icosaedro.
La incidencia no se deja verbal: se fija por
\begin{equation}
\label{rm:eq:gravity-icosahedron-incidence}
\mathcal I_{12}^{\rm ico}
=\{(v,w)\in\Omega_{12}^{\rm ico}\times\Omega_{12}^{\rm ico}:
v\neq w,\ \langle v,w\rangle=1/\sqrt5\}.
\end{equation}
Cada v\'ertice tiene cinco vecinos en esta relaci\'on. El puente de lectura
se tipa como una aplicaci\'on
\begin{equation}
\label{rm:eq:gravity-exceptional-reader-map}
\pi_{\rm ico}:X_{\rm TPK}^{\rm exc}\longrightarrow\Omega_{12}^{\rm ico}.
\end{equation}
La cardinalidad doce de la imagen de \(\pi_{\rm ico}\) no es suficiente.
Para definir una realizaci\'on de incidencia, la aplicaci\'on debe satisfacer
simult\'aneamente:
\begin{enumerate}[label=(\roman*)]
\item ser sobreyectiva y constante exactamente en las fibras que descartan
la carta aritm\'etica pero conservan la ruta excepcional;
\item entrelazar la acci\'on de los automorfismos orientados del antecedente
con una acci\'on transitiva sobre \(\Omega_{12}^{\rm ico}\);
\item transportar la incidencia TPK declarada a
\(\mathcal I_{12}^{\rm ico}\), no
s\'olo sus cardinales;
\item conservar la involuci\'on de hoja y el retorno non\'adico en el
estabilizador de la fibra correspondiente.
\end{enumerate}

Estas condiciones separan tres afirmaciones: la existencia del mapa
\(\pi_{\rm ico}\), la identificaci\'on del grupo que act\'ua y la
representaci\'on lineal que se construye despu\'es. La segunda y la tercera
no pueden usarse para definir retrospectivamente la primera.

\begin{theorem}[Descenso de incidencia al cociente icosa\'edrico]
\label{rm:thm:gravity-exceptional-descent}
Sup\'ongase que \(\pi_{\rm ico}\) satisface (i)--(iv) y que la acci\'on
orientada inducida en \(\Omega_{12}^{\rm ico}\) tiene grupo
\(G_{\rm ico}\simeq A_5\). Entonces, para toda clase
\(v_0\in\Omega_{12}^{\rm ico}\),
\begin{equation}
\label{rm:eq:gravity-icosahedral-quotient}
H_{v_0}=\operatorname{Stab}_{G_{\rm ico}}(v_0)\simeq C_5,
\qquad
\Omega_{12}^{\rm ico}\simeq G_{\rm ico}/H_{v_0}\simeq A_5/C_5.
\end{equation}
La aplicaci\'on
\begin{equation}
\label{rm:eq:gravity-coset-map}
\kappa_{v_0}:G_{\rm ico}/H_{v_0}\longrightarrow\Omega_{12}^{\rm ico},
\qquad gH_{v_0}\longmapsto g\cdot v_0,
\end{equation}
es una biyecci\'on equivariante y preserva la incidencia si en el cociente
se declara
\begin{equation}
\label{rm:eq:gravity-coset-incidence}
gH_{v_0}\sim g'H_{v_0}
\quad\Longleftrightarrow\quad
(g\cdot v_0,g'\cdot v_0)\in\mathcal I_{12}^{\rm ico}.
\end{equation}
Por tanto, el compuesto \(\kappa_{v_0}^{-1}\pi_{\rm ico}\) es el mapa
tipado del antecedente TPK al cociente \(A_5/C_5\).
\end{theorem}

\begin{proof}
La lista de \cref{rm:eq:gravity-icosahedron-vertices} tiene doce elementos.
Por transitividad, \(|G_{\rm ico}\cdot v_0|=12\), y el teorema
\`orbita--estabilizador da \(|H_{v_0}|=60/12=5\). Todo grupo de orden
primo cinco es c\'iclico; por tanto, \(H_{v_0}\simeq C_5\). La f\'ormula de
\(\kappa_{v_0}\) no depende del representante: si
\(gH_{v_0}=g'H_{v_0}\), entonces \(g^{-1}g'\in H_{v_0}\) y
\(g'v_0=gv_0\). La sobreyectividad procede de la transitividad, y la
igualdad de cardinales da inyectividad. La equivarianza es inmediata:
\(\kappa_{v_0}(agH_{v_0})=a\kappa_{v_0}(gH_{v_0})\). Finalmente,
\cref{rm:eq:gravity-coset-incidence} es precisamente el transporte de
estructura por esa biyecci\'on; la hip\'otesis (iii) hace conmutar este
transporte con \(\pi_{\rm ico}\).
\end{proof}

El teorema no sustituye la construcci\'on de \(\pi_{\rm ico}\): demuestra
qu\'e se obtiene una vez verificados sus cuatro requisitos. En el perfil
HMT completo, el corredor de incidencia que llega a
Paley--Witt--Golay--Leech--\(V^\natural\)--Monstruo--Moonshine es el
dominio de definici\'on de ese mapa. En este volumen aut\'onomo, la entrada exacta es
el mapa \cref{rm:eq:gravity-exceptional-reader-map} con sus condiciones; todo
el tramo desde \(A_5/C_5\) hasta la reducci\'on TT se prueba a continuaci\'on
sin volver a consultar un valor gravitatorio.

La cadena de procedencia se expresa mediante todos sus morfismos:
\begin{equation}
\label{rm:eq:gravity-exceptional-provenance}
\APP\longrightarrow\TRIT\longrightarrow\TPK
\longrightarrow X_{\rm TPK}^{\rm exc}
\xrightarrow{\ \pi_{\rm ico}\ }\Omega_{12}^{\rm ico}
\xleftarrow[\simeq]{\ \kappa_{v_0}\ }A_5/C_5.
\end{equation}
El corredor excepcional completo se demuestra en su residencia propia. La
hip\'otesis que este cap\'itulo no puede reemplazar por un conteo es la
existencia de \(\pi_{\rm ico}\) con (i)--(iv). Fijada esa entrada, el descenso
al cociente es el \cref{rm:thm:gravity-exceptional-descent}, y se construyen
desde \`el \(P_5\), \(\Gamma_5\) y \(\Lambda(n)\). La realizaci\'on del
portador resultante como campo gravitatorio sigue siendo una interfaz
f\'isica tipada, no una consecuencia de cardinalidad.

Sea \(G=A_5\), sea \(H\simeq C_5\) el estabilizador de un v\'ertice
orientado del icosaedro y consid\'erese
\begin{equation}
\label{rm:eq:gravity-v12}
V_{12}=\R[G/H],
\qquad |G/H|=\frac{60}{5}=12.
\end{equation}
La representaci\'on de permutaci\'on \(\rho_{12}\) se descompone como
\begin{equation}
\label{rm:eq:gravity-v12-decomposition}
V_{12}\simeq\mathbf1\oplus V_3\oplus V_{3'}\oplus V_5.
\end{equation}
El idempotente central del sumando de car\'acter
\(\chi_5=(5,1,-1,0,0)\) es
\begin{equation}
\label{rm:eq:gravity-p5}
\boxed{
P_5=\frac5{60}\sum_{g\in A_5}\chi_5(g^{-1})\rho_{12}(g)
=\frac1{12}\left(
5I+\sum_{g\in2A}\rho_{12}(g)
-\sum_{g\in3A}\rho_{12}(g)
\right).}
\end{equation}

\begin{theorem}[Proyector y acoplamiento gravitatorio de rango cinco]
\label{rm:thm:gravity-p5-coupling}
El operador \(P_5\) satisface
\begin{equation}
\label{rm:eq:gravity-p5-properties}
P_5^2=P_5=P_5^*,
\qquad \Tr P_5=\rank P_5=5.
\end{equation}
Para
\begin{equation}
\label{rm:eq:gravity-G5}
\mathbb G_{5,a}:=G_a(\theta_{108}^{\rm circ})P_5
\end{equation}
se tiene
\begin{equation}
\label{rm:eq:gravity-G5-spectrum}
\Spec(\mathbb G_{5,a})
=\{G_a(\theta_{108}^{\rm circ})\text{ con multiplicidad }5,
\;0\text{ con multiplicidad }7\}.
\end{equation}
\end{theorem}

\begin{proof}
La ortogonalidad de caracteres convierte \cref{rm:eq:gravity-p5} en el
idempotente central que act\'ua como la identidad sobre la \`unica copia de
\(V_5\) y como cero en los otros tres sumandos de
\cref{rm:eq:gravity-v12-decomposition}. La realidad del car\'acter y la
ortogonalidad de \(\rho_{12}\) dan la autoadjunci\'on. Su traza y su rango
son cinco. Multiplicar un proyector de rango cinco por el escalar positivo
\(G_a\) produce inmediatamente el espectro indicado.
\end{proof}

La expresi\'on \(\mathbb G_{5,a}\) no a\~nade cinco valores diferentes de
la constante gravitatoria. Es una misma secci\'on de \(G\) actuando sobre
el sumando tensorial seleccionado, y anulando los otros siete modos del
m\'odulo de permutaci\'on.
